\documentclass[11pt,a4paper,sans]{moderncv} 
\moderncvstyle{classic} 
\moderncvcolor{blue}
\AtBeginDocument{\hypersetup{hidelinks}}

\usepackage[utf8]{inputenc} 
\usepackage[scale=0.80]{geometry} 
\usepackage{xcolor}


\name{Ian Jordy Lopez Diaz}{}
\title{Cientista de Dados}
\phone[mobile]{+55~(48)~99169~9996} 
\email{ianlopezdiaz@gmail.com}
\social[linkedin]{ianlopezdiaz}
\social[github]{ianlopezdiaz}
\homepage{ianlopezdiaz.github.io/portfolio}
%\social[linkedin]{linkedin.com/in/ianlopezdiaz}
%\social[github]{github.com/ianlopezdiaz}

\begin{document} 

\makecvtitle 

%-----------------------------------
% Resumo Profissional
%-----------------------------------
\section{Resumo Profissional}
\cvitem{}{
Cientista de Dados com doutorado em Física (Mecânica Estatística) e sólida experiência na aplicação de machine learning,
modelagem estatística e métodos numéricos a problemas industriais reais.
Atualmente atuo na AQTech Power Prognostics, desenvolvendo soluções orientadas a dados para o setor de energia,
incluindo modelagem preditiva, análise de sinais e construção de pipelines de dados escaláveis.
}

%-----------------------------------
% Habilidades Técnicas
%-----------------------------------
\section{Habilidades Técnicas}
\cvitem{Programação}{Python, SQL, R, C/C++}
\cvitem{Ciência de Dados}{Pandas, NumPy, SciPy, Scikit-Learn, TensorFlow, PyTorch, Statsmodels, Pipelines de ETL, Limpeza de Dados, Engenharia de Atributos}
\cvitem{Bancos de Dados}{PostgreSQL, MySQL, SQLite}
\cvitem{Visualização}{Matplotlib, Plotly, Seaborn}
\cvitem{Ferramentas}{Git, Linux, Jupyter, Docker}
\cvitem{Computação Científica}{Simulações de Monte Carlo, Métodos Numéricos, Modelagem Estatística}

%-----------------------------------
% Experiência
%-----------------------------------
\section{Experiência}

\cventry{2026--Presente}{Cientista de Dados}{AQTech Power Prognostics}{Florianópolis, Brasil}{}{
\begin{itemize}
\item Desenvolvimento de modelos de machine learning para manutenção preditiva e detecção de falhas em sistemas de energia.
\end{itemize}
}

\cventry{2021--2024}{Cientista/Analista de Dados/Desenvolvedor}{Lima Brain Technology}{Florianópolis, Brasil}{}{
\begin{itemize}
\item Desenvolvimento e implementação de pipelines de dados, e modelos de previsão.
\item Desenvolvimento de simulações de Monte Carlo e algoritmos numéricos.
\end{itemize}}

\cventry{2008--2025}{Professor de Física}{}{}{}{
\begin{itemize}
\item UFSC, Florianópolis, Brasil; UDESC, Ibirama, Brasil; UFFS, Cerro Largo, Brasil
%\item UFSC, Florianópolis, Brazil:
%08/2008--07/2010; 02/2013--12/2014; 02/2018--12/2019; 08/2023--07/2025.
%\item UDESC, Ibirama, Brazil:
%02/2020--07/2022.
%\item UFFS, Cerro Largo, Brazil:
%01/2011--02/2013.
\end{itemize}}



%-----------------------------------
% Educação
%-----------------------------------
\section{Educação} 
\cventry{2015--2017}{Doutorado, Física (Mecânica Estatística)}{UFSC}{Florianópolis, SC}{}{} 
\cventry{2007--2009}{Mestrado, Física (Mecânica Estatística)}{UFSC}{Florianópolis, SC}{}{} 
\cventry{2003--2006}{Bacharelado, Física}{UFSC}{Florianópolis, SC}{}{} 

%-----------------------------------
% Publicações
%-----------------------------------
\section{Publicações com Revisão de Pares}
%{\small
Autor de 5 artigos revisados por pares em Física Estatística e simulações Monte Carlo, publicados em periódicos como \textit{Physica A}, \textit{Physica B} e \textit{Physical Review E}.
%}



\end{document}

